% Part: first-order-logic
% Chapter: axiomatic-deduction
% Section: axioms-rules-quantifiers

\documentclass[../../../include/open-logic-section]{subfiles}

\begin{document}

\olfileid[fr]{fol}{axd}{qua}

\olsection{Axiomes et règles pour les quantificateurs}

\begin{defn}[Axiomes des quantificateurs]
Les \emph{axiomes} qui régissent les quantificateurs sont toutes les
instances des schémas suivants :
\begin{align}
\ollabel{ax:q1} & \lforall[x][!B] \lif !B(t), \\
\ollabel{ax:q2} & !B(t) \lif \lexists[x][!B].
\end{align}
où $t$ est un terme clos quelconque.
\end{defn}

\begin{defn}[Règles des quantificateurs]
\begin{enumerate}
\item Si $!B \lif !A(a)$ figure déjà dans la !!{derivation} et si $a$
  ne figure ni dans~$\Gamma$ ni dans~$!B$, alors
  $!B \lif \lforall[x][!A(x)]$ est une étape d'inférence correcte.
\item Si $!A(a) \lif !B$ figure déjà dans la !!{derivation} et si $a$
  ne figure ni dans~$\Gamma$ ni dans~$!B$, alors
  $\lexists[x][!A(x)] \lif !B$ est une étape d'inférence correcte.
\end{enumerate}
\end{defn}

Nous abrégerons chacune de ces deux règles par « \QR ».

\begin{editorial}
Dans la source, les deux clauses de la définition précédente sont des
commandes \texttt{\string\item} placées hors de tout environnement de
liste. L'énumération attendue a été rétablie; les deux règles et leurs
conditions de fraîcheur sont inchangées.
\end{editorial}

\end{document}
